\scanpage{131}
\mspage{80}

de la remarque \deleted{\unsure{doit}} signifie simplement que
le dérivateur dual $\mathrm{HOT}^{\circ}$ ne possède
pas cette même propriété. Je dois \unsure{envisager|renvoyer} à ce
sujet \deleted{\unsure{aux}} \deleted{\unsure{l'}}\inserted{deux} axiomes pour les dérivateurs correspondants à $\mathrm W(7\ \mathrm{bis})$, $\mathrm W(7)$, ce que je ferai
sans doute sous peu, de façon à donner
un sens plus précis au présent commentaire.

2.2. \uline{Application aux carrés cocartésiens déduits
d'une immersion ouverte ou fermée.}

Les axiomes $\mathrm W(7)$, $\mathrm W(7\ \mathrm{bis})$ (du \unsure{cube|crible}, \deleted{\unsure{carré}} et du
carré cocartésien) concernent des carrés cocartésiens
de nature très particulière. Nous allons
voir à présent comment en déduire des
énoncés \unsure{de|du} \unsure{=|même} type pour des carrés cocartésiens
beaucoup plus généraux, y compris pour des
carrés qui n'ont rien de cocartésien \inserted{(dans $\mathrm{Cat}$)} --
\deleted{\unsure{mais}} \inserted{et} que nous appellerons ``homotopiquement
cocartésiens'' (relativement à $\mathcal W$), ou simplement
\unsure{$\mathcal W$-cocartésiens}. (Il y aura également des carrés
$\mathcal W$-cartésiens un jour\ldots) Dans le cas où
$\mathcal W=\mathcal W_\infty$, je les appellerai \inserted{aussi} $\mathrm{Hot}$-cocartésiens.

Je commencerai par le cas particulier
d'un carré cocartésien associé à une \unsure{imm.|immersion}
ouverte, (diagramme \unsure{(5)|(*)|(6)} p.~74) :

\diagram{\(\text{(\unsure{*|**})}\qquad\begin{tikzcd}[ampersand replacement=\&]
U \arrow[r,hook,"i"] \arrow[d,"p"'] \& X \arrow[d,"q"] \& T \arrow[l,hook',dotted,"j"'] \arrow[d,dotted] \\
U' \arrow[r,hook,"i'"] \& X' \& T'\unsure{\simeq|=}T \arrow[l,hook',dotted,"j'"']
\end{tikzcd}\)}.

% archive sheet number: 131
